% !TEX program = pdflatex
% TD - Les suites (Première, tronc commun) -- SQUELETTE
% Compilation : pdflatex TD-suites.tex
\documentclass[11pt,a4paper]{article}

\usepackage[utf8]{inputenc}
\usepackage[T1]{fontenc}
\usepackage[french]{babel}
\usepackage{lmodern}
\usepackage{amsmath,amssymb,amsfonts}
\usepackage[a4paper,margin=2.2cm]{geometry}
\usepackage{enumitem}
\usepackage{fancyhdr}
\usepackage{xcolor}
\usepackage{array}
\usepackage{graphicx}       % pour insérer les captures
% \graphicspath{{images/}}  % décommenter si les captures sont dans un sous-dossier

%------------------------------------------------------------
% Mise en forme
%------------------------------------------------------------
\definecolor{bleu}{RGB}{30,80,140}

\pagestyle{fancy}
\fancyhf{}
\lhead{\small Mathématiques -- Première}
\rhead{\small TD1 : Les suites}
\cfoot{\small \thepage}
\renewcommand{\headrulewidth}{0.4pt}

% Environnement "exercice" numéroté avec titre
\newcounter{exo}
\newcommand{\exo}[1]{%
  \refstepcounter{exo}%
  \vspace{0.8cm}
  \noindent{\color{bleu}\large\bfseries Exercice \theexo\ -- #1}\par
  \nopagebreak\smallskip
}

% Raccourci pour insérer une capture (légende optionnelle en 2e argument)
% Utilisation :  \capture{nom-fichier.png}{0.6}
\newcommand{\capture}[2]{%
  \begin{center}\includegraphics[width=#2\textwidth]{#1}\end{center}}

\setlength{\parindent}{0pt}
\setlist[enumerate]{leftmargin=*,itemsep=4pt,topsep=4pt}

%------------------------------------------------------------
\begin{document}

%------------------- En-tête -------------------
\begin{center}
  {\color{bleu}\LARGE\bfseries TD -- Les suites}\\[4pt]
  {\large Première -- tronc commun}
\end{center}

\vspace{6pt}
\noindent\fbox{\parbox{\dimexpr\textwidth-2\fboxsep-2\fboxrule}{%
  \small \textbf{Objectifs du TD :} Comprendre les bases des suites et les cas particuliers
  des suites arithmétiques.}}

%============================================================
%  EXERCICE 1
%============================================================
\exo{Représentation graphique d'une suite}
% --- Énoncé ---
\capture{exo1.png}{0.5}   % <-- insérer une capture ici
\begin{enumerate}
  \item Parmis les représentations graphiques suivantes, lesquelles représentent
  des modèles discrets à croissance linéaire ? Justifier.
\end{enumerate}

%============================================================
%  EXERCICE 2
%============================================================
\exo{Calculs de termes}
% --- Énoncé ---
Pour chacune de ces suites arithmétiques, calculer $u_{10}$.
\begin{enumerate}
  \item $u_n = 5n - 12$
  \item $\begin{cases} u_0 = 2 \\ u_{n+1} = u_n + 4 \end{cases}$
  \item La raison vaut $3$ et $u_8 = -4$
\end{enumerate}

%============================================================
%  EXERCICE 3
%============================================================
\exo{Raison et variations d'une suite arithmétique}
% --- Énoncé ---
% \capture{capture-2.png}{0.6}
Précuser la raison de chacune des suites arithmétiques suivantes dont on donne les premiers
termes.

Indiquer le sens de variation de cette suite.
\begin{enumerate}
  \item $1$ ; $-1$ ; $-3$ ; ...
  \item $\dfrac{3}{5}$ ; $1$ ; $\dfrac{7}{5}$ ; $\dfrac{9}{5}$ ; ...
\end{enumerate}

%============================================================
%  EXERCICE 4
%============================================================
\exo{Modélisation}
% --- Énoncé ---

La réserve de noisettes d'un écureuil diminue chaque jour de l'hiver.

\vspace{5pt}
Elle débute à 156 noisettes. Chaque mois, la réserve diminue de 23 noisettes.

\begin{enumerate}
  \item Peut-on modéliser cette situation par une suite arithmétique ?
  \item Donner les éléments caractéristiques de cette suite.
  \item La réserve tiendra-t-elle l'hiver ?
\end{enumerate}

%============================================================
%  (dupliquer un bloc \exo{...} pour ajouter un exercice)
%============================================================

\end{document}
